% GENERATED FILE. Edit canonical lessons, then run npm run generate:books.
% canonical-source-hash: fnv1a64:e2aeb0f85f3e44f7
% canonical-lessons: MW-C65-kivar, MW-C65-rasoi, MW-C65-kamro, MW-C65-kitab, MW-C65-kalam

\chapter{Door, Kitchen, Room, Book, Pen}
\label{ch:mw-door-kitchen-room-book-pen}
\hlchaptermodality{65}

\begin{chapteropening}
\textbf{By the end of this chapter:} \emph{I can say a door, a kitchen, a room, a book and a pen.}

Complete the last lesson of chapter 65: I can say a door, a kitchen, a room, a book and a pen.
\end{chapteropening}

\section[\texorpdfstring{\mw{किवाड़} (kivāṛ)}{किवाड़ (kivāṛ)}]{\mw{किवाड़} (kivāṛ) — a door}

\label{lesson:MW-C65-kivar}



\begin{quote}
Before the new one: say the Marwadi for a knife, then the Marwadi for a bottle.
\end{quote}

\subsection*{You'll want to know: \mw{किवाड़}}
\textbf{\mw{किवाड़}} — \emph{kivāṛ} — \textquotedblleft{}a door\textquotedblright{}.

\begin{cousinweb}
\textbf{\mw{किवाड़}} (\emph{kivāṛ}) is a tadbhava, a word inherited from Sanskrit and reshaped by sound change. It comes from \emph{kapāṭa}, \textquotedblleft{}leaf of a door\textquotedblright{}. CDIAL lists Old Marwari \emph{kibāṛi} under it (CDIAL 2963). CDIAL notes a possible Dravidian origin for the Sanskrit word.
\end{cousinweb}

\begin{grammarlens}[title={What you've built}]
One more word for this chapter.
\end{grammarlens}

\subsection*{Your turn}
\begin{itemize}
\raggedright
  \item \emph{Say it:} \emph{kivāṛ}
  \item \emph{Say it:} \emph{kivāṛ}, once more
  \item \emph{Say it:} say \emph{kivāṛ}
  \item \emph{Recall:} say the Marwadi for a knife, then the Marwadi for a bottle, then say \emph{kivāṛ} again
\end{itemize}

\begin{tcolorbox}[breakable,colback=teal!4,colframe=teal!35!black,title={Before you move on}]
What does \mw{किवाड़} mean? (\textquotedblleft{}A door\textquotedblright{}.) Say it once more.
\end{tcolorbox}
\section[\texorpdfstring{\mw{रसोई} (rasoī)}{रसोई (rasoī)}]{\mw{रसोई} (rasoī) — a kitchen}

\label{lesson:MW-C65-rasoi}



\begin{quote}
Before the new one: say the Marwadi for a bottle, then the Marwadi for a door.
\end{quote}

\subsection*{You'll want to know: \mw{रसोई}}
\textbf{\mw{रसोई}} — \emph{rasoī} — \textquotedblleft{}a kitchen\textquotedblright{}.

\begin{cousinweb}
\textbf{\mw{रसोई}} (\emph{rasoī}) is a tadbhava, a word inherited from Sanskrit and reshaped by sound change. It comes from \emph{rasavatī}, \textquotedblleft{}kitchen\textquotedblright{}. CDIAL lists Marwari \emph{rasoī} under it (CDIAL 10656). It is built on \emph{rasa}, \textquotedblleft{}juice, flavour\textquotedblright{}.
\end{cousinweb}

\begin{grammarlens}[title={What you've built}]
One more word for this chapter.
\end{grammarlens}

\subsection*{Your turn}
\begin{itemize}
\raggedright
  \item \emph{Say it:} \emph{rasoī}
  \item \emph{Say it:} \emph{rasoī}, once more
  \item \emph{Say it:} say \emph{rasoī}
  \item \emph{Recall:} say the Marwadi for a bottle, then the Marwadi for a door, then say \emph{rasoī} again
\end{itemize}

\begin{tcolorbox}[breakable,colback=teal!4,colframe=teal!35!black,title={Before you move on}]
What does \mw{रसोई} mean? (\textquotedblleft{}A kitchen\textquotedblright{}.) Say it once more.
\end{tcolorbox}
\section[\texorpdfstring{\mw{कमरो} (kamro)}{कमरो (kamro)}]{\mw{कमरो} (kamro) — a room}

\label{lesson:MW-C65-kamro}



\begin{quote}
Before the new one: say the Marwadi for a door, then the Marwadi for a kitchen.
\end{quote}

\subsection*{You'll want to know: \mw{कमरो}}
\textbf{\mw{कमरो}} — \emph{kamro} — \textquotedblleft{}a room\textquotedblright{}.

\begin{cousinweb}
\textbf{\mw{कमरो}} (\emph{kamro}) is a loan from Portuguese \emph{câmara}, \textquotedblleft{}room, chamber\textquotedblright{}. It goes back to Latin \emph{camera}, \textquotedblleft{}vaulted room\textquotedblright{}.
\end{cousinweb}

\begin{grammarlens}[title={What you've built}]
One more word for this chapter.
\end{grammarlens}

\subsection*{Your turn}
\begin{itemize}
\raggedright
  \item \emph{Say it:} \emph{kamro}
  \item \emph{Say it:} \emph{kamro}, once more
  \item \emph{Say it:} say \emph{kamro}
  \item \emph{Recall:} say the Marwadi for a door, then the Marwadi for a kitchen, then say \emph{kamro} again
\end{itemize}

\begin{tcolorbox}[breakable,colback=teal!4,colframe=teal!35!black,title={Before you move on}]
What does \mw{कमरो} mean? (\textquotedblleft{}A room\textquotedblright{}.) Say it once more.
\end{tcolorbox}
\section[\texorpdfstring{\mw{किताब} (kitāb)}{किताब (kitāb)}]{\mw{किताब} (kitāb) — a book}

\label{lesson:MW-C65-kitab}



\begin{quote}
Before the new one: say the Marwadi for a kitchen, then the Marwadi for a room.
\end{quote}

\subsection*{You'll want to know: \mw{किताब}}
\textbf{\mw{किताब}} — \emph{kitāb} — \textquotedblleft{}a book\textquotedblright{}.

\begin{cousinweb}
\textbf{\mw{किताब}} (\emph{kitāb}) is a loan from Arabic \emph{kitāb}, \textquotedblleft{}book\textquotedblright{}, which reached Marwadi through Persian. Its Arabic root is \emph{k-t-b} (Lane).
\end{cousinweb}

\begin{grammarlens}[title={What you've built}]
One more word for this chapter.
\end{grammarlens}

\subsection*{Your turn}
\begin{itemize}
\raggedright
  \item \emph{Say it:} \emph{kitāb}
  \item \emph{Say it:} \emph{kitāb}, once more
  \item \emph{Say it:} say \emph{kitāb}
  \item \emph{Recall:} say the Marwadi for a kitchen, then the Marwadi for a room, then say \emph{kitāb} again
\end{itemize}

\begin{tcolorbox}[breakable,colback=teal!4,colframe=teal!35!black,title={Before you move on}]
What does \mw{किताब} mean? (\textquotedblleft{}A book\textquotedblright{}.) Say it once more.
\end{tcolorbox}
\section[\texorpdfstring{\mw{कलम} (kalam)}{कलम (kalam)}]{\mw{कलम} (kalam) — a pen}

\label{lesson:MW-C65-kalam}



\begin{quote}
Before the new one: say the Marwadi for a room, then the Marwadi for a book.
\end{quote}

\subsection*{You'll want to know: \mw{कलम}}
\textbf{\mw{कलम}} — \emph{kalam} — \textquotedblleft{}a pen\textquotedblright{}.

\begin{cousinweb}
\textbf{\mw{कलम}} (\emph{kalam}) is a loan from Arabic \emph{qalam}, \textquotedblleft{}reed pen\textquotedblright{}, which reached Marwadi through Persian. Its Arabic root is \emph{q-l-m} (Lane). CDIAL marks Hindi \emph{kalam}, \textquotedblleft{}pen\textquotedblright{}, as a loan (CDIAL 2653).
\end{cousinweb}

\begin{grammarlens}[title={What you've built}]
One more word for this chapter.
\end{grammarlens}

\subsection*{Your turn}
\begin{itemize}
\raggedright
  \item \emph{Say it:} \emph{kalam}
  \item \emph{Say it:} \emph{kalam}, once more
  \item \emph{Say it:} say \emph{kalam}
  \item \emph{Recall:} say the Marwadi for a room, then the Marwadi for a book, then say \emph{kalam} again
\end{itemize}

\begin{tcolorbox}[breakable,colback=teal!4,colframe=teal!35!black,title={Before you move on}]
What does \mw{कलम} mean? (\textquotedblleft{}A pen\textquotedblright{}.) Say it once more.
\end{tcolorbox}
